\phantomsection
\subsection*{M2E. Other Fruit Mead}
\addcontentsline{toc}{subsection}{M2E. Other Fruit Mead}

\textit{O estilo Other Fruit Mead é destinado a meloméis elaborados com qualquer fruta não associada a nenhum estilo anterior de melomel, ou com uma combinação de frutas de múltiplos estilos de hidromel com frutas (como uvas e frutas de caroço). Alguns exemplos incluem frutas cítricas, frutas secas (tâmara, ameixa seca, passa, etc.), figo, romã, figo-da-índia, banana, abacaxi, goiaba e a maioria das demais frutas tropicais. Em caso de dúvida, inscreva a fruta nesta categoria — os juízes devem ser flexíveis ao permitir frutas não explicitamente mencionadas em outros estilos.}

\textbf{Impressão Geral}: Como um M1 Traditional Mead com um caráter frutado perceptível. Uma ampla gama de resultados é possível, dependendo das escolhas do hidromel base e dos ingredientes adicionais, mas o produto final deve ser equilibrado, agradável e prazeroso. O caráter da fruta não deve dominar completamente o hidromel.

\textbf{Aroma}: Os aromas da fruta podem ser percebidos como desde um suco doce e fresco até um vinho de fruta envelhecido ou do Porto, dependendo da força alcoólica e do dulçor. Se variedades de fruta forem declaradas, o caráter varietal deve estar perceptível. Se múltiplos tipos de fruta forem declarados, eles não precisam estar em proporção igual, tendo em vista que algumas frutas são mais fortes e distintivas que outras.

O caráter do mel pode variar de sutil a intenso, dependendo da força e dulçor, podendo expressar o néctar floral típico de eventuais variedades de mel declaradas. O buquê de fermentação deve ser equilibrado. Hidroméis mais fortes podem apresentar uma leve nota alcoólica picante, e os carbonatados tendem a expressar mais aroma. Os aromas provenientes da fruta devem complementar e realçar o hidromel base, não sobrepujá-lo.

\textbf{Aparência}: Limpidez de boa a brilhante; quanto maior a limpidez, mais desejável. A presença de bolhas, espuma, colarinho ou efervescência depende do nível de carbonatação declarado. A cor é derivada da variedade de fruta e de mel utilizadas. Quanto maior a vivacidade e o brilho da cor, mais desejável.

\textbf{Sabor}: Semelhante ao Aroma quanto a característica e intensidade de fruta e mel, se tornando mais proeminente em versões mais alcoólicas e doces. O caráter da fruta deve ser perceptível e equilibrado com o nível de dulçor do hidromel. Algumas frutas podem conferir gostos básicos (amargo, doce, ácido) além de suas características típicas. Muitas frutas conferem acidez e taninos naturais que precisam ser adequadamente equilibrados com o hidromel base.

Dulçor residual e final de acordo com o nível de dulçor declarado. A acidez deve ser equilibrada, macia ou vibrante, mas não agressiva. Taninos podem fazer um hidromel suave parecer mais seco. Características ásperas, desagradáveis ou excessivamente sulfurosas ou com notas de levedura/autólise são indesejáveis.

\textbf{Sensação na Boca}: A carbonatação é indicada pelo nível declarado, de tranquilo a espumante. O corpo geralmente aumenta com o dulçor e a força alcoólica, tipicamente de médio-leve a cheio. A percepção de álcool acompanha a força alcoólica declarada, variando de ausente a perceptível e com sensação de aquecimento. Os taninos de frutas mais escuras podem conferir certo corpo e uma sensação de secura, mas isso não deve ser excessivo.

\textbf{Ingredientes}: Os mesmos ingredientes de um M1 Traditional Mead, com a adição de pelo menos um tipo de fruta que não seja fruta pomácea, uva, baga ou fruta de caroço.

\textbf{Comentários}: Este é um estilo abrangente para a categoria Melomel, portanto apenas meloméis que não se encaixem em nenhum estilo anterior de hidromel frutado devem ser inscritos aqui. Versões com especiarias devem ser inscritas como M3C Fruit and Spice Mead. Hidroméis com frutas e outros ingredientes que não forem frutas provavelmente devem ser inscritos na categoria M4F (Experimental Mead).

\textbf{Instruções para Inscrição}: Os participantes devem especificar os níveis de dulçor, carbonatação e força alcoólica. Variedades de mel podem ser declaradas. A fruta deve obrigatoriamente ser declarada.

\textbf{Exemplos Comerciais}: Intermiel Black Currant and Citrus Sparkling Mead, Lost Cause Malum, Moonlight Desire, Redstone Blood Orange Nectar, Redstone Passion Fruit Nectar, Schramm's Mead Apple Crandall Reserve.
